\section{Conclusion}
This study separates historical training observations from search-time improvements in VGG16--GeM retrieval. The validation-selected B0 remains a strong baseline, while the two new training runs exhibit late validation deterioration. On the completed B0 corruption study, AQE raises dataset-macro mAP from 77.20\% to 80.64\%. Sparse diffusion reaches 83.31\%, with a strong Paris gain but a clean Oxford loss. The same method therefore benefits some encoder/dataset combinations more than others.

An independent implementation audit identifies dense query seeding as the cause of the severe archived diffusion collapse. Correct seed support recovers useful rankings, yet it does not make diffusion uniformly beneficial. Archived local/geometric experiments also show that feedback quality and query cost matter: verified expansion improves the weaker epoch-100 encoders but costs seconds per query. Validation-based checkpoint selection, faithful implementation and separate quality/latency accounting are necessary for useful comparisons. The released source, numerical records and solver checks make these findings traceable.
